\documentclass[10pt,a4paper]{amsart}
\usepackage[margin=19mm]{geometry}
\usepackage{amsmath,amssymb,mathtools}
\usepackage[scr=rsfs,cal=euler]{mathalfa}
\usepackage{xfrac}
\usepackage[unicode,hidelinks,bookmarks=false]{hyperref}
\newcommand{\ie}{i.e.\ }
\newcommand{\cf}{cf.\ }
\newcommand{\resp}{respectively\ }
\newcommand{\Subsection}{\subsection}
\providecommand{\nobreakdash}{\nobreak\hbox{-}}
\setlength{\emergencystretch}{2em}
\multlinegap0pt
\theoremstyle{plain}
\newtheorem{thm}{Theorem}
\newtheorem{prop}{Proposition}
\newtheorem{lem}{Lemma}
\newtheorem{cor}{Corollary}
\newtheorem{claim}{Claim}
\newtheorem{conj}{Conjecture}
\theoremstyle{remark}
\newtheorem{rem}{Remark}
\title[Random burning of the Euclidean torus]{Random burning of the Euclidean torus: the coupon-collector description of the burning time}
\author{Guillaume Blanc, Alice Contat}
\date{Source: arXiv:2509.02562v2 (CC BY 4.0)}
\usepackage{stmaryrd}
\begin{document}
\maketitle
\noindent\emph{Source and licence.} Random burning of the Euclidean torus. Version arXiv:2509.02562v2 (CC BY 4.0), distributed under the Creative Commons Attribution 4.0 International licence, \url{http://creativecommons.org/licenses/by/4.0/}, as recorded in the arXiv licence metadata for this exact version and shown on the abstract page \url{https://arxiv.org/abs/2509.02562v2}. Full licence notice: licenses/arxiv-2509.02562-CC-BY-4.0.txt.\par\smallskip

\noindent\textit{Editorial scope.} Counted unit: the article's Theorem with source label \texttt{thm:couponcollector} (source0.tex lines 115--118), the coupon-collector description of the burning time of the Euclidean torus, delivered together with the article's complete original proof of it (source0.tex lines 296--414, the article's own proof section whose heading names the counted theorem; the article places this proof after the auxiliary results it uses). No mathematics is written, repaired or re-derived: statement and proof are the article's own text line for line, in the article's own notation (the burning process, the covered and uncovered regions, the coupon-collector times, the limit statements), and no retained line is substituted, re-flowed or omitted. Required context retained uncounted: the statement of the article's deterministic bound on the covered region (lines 76--81), which the counted proof invokes. The proof of that context result is not delivered. Nothing else of the article is retained: its other results, its simulations and its figures are omitted. No citation command occurs in any retained line, so the article's bibliography is not delivered and the wrapper carries no bibliography entry, although the article ships one. Every cross-reference inside the retained spans resolves inside the delivered document. Editorial formatting only: an AMS \texttt{amsart} preamble that restates the environments the retained lines use, namely the article's declarations of Theorem, Proposition, Lemma, Corollary, Claim, Conjecture and Remark, and the macros that the retained lines need. The article is typeset in its own \texttt{article}-class format with its own fonts; no class or style file is shipped with this package and the wrapper is an amsart document. Numbering differs from the article, because the article numbers its theorem-like environments with independent counters and only a subset is retained: the retained context proposition prints as Proposition 1 and the counted theorem as Theorem 1. No picture environment and no graphical inclusion occurs in any retained line, and the package delivers no image asset. One bundled source file, \texttt{sources/184337/source0.tex}, accompanies the delivered item as the exact source of every retained span. Every retained span is recorded in \texttt{delivery-evidence.json} with method \texttt{exact\_tex}.
\begin{prop}\label{prop:detbounds}
As $n\to\infty$, we have
\[\left(\alpha_d+o(1)\right)\cdot n^{d/(d+1)}\leq b\left(\mathbb{T}_n^d\right)\leq\left(\gamma_d+o(1)\right)\cdot n^{d/(d+1)},\]
where
\[\alpha_d=\left(\frac{(d+1)!}{2^d}\right)^{1/(d+1)}\quad\text{and}\quad\gamma_d=\left(\left(1+\frac{1}{d}\right)^d\cdot(d+1)!\right)^{1/(d+1)}.\]
\end{prop}

\begin{thm}\label{thm:couponcollector}
As $n\to\infty$, we have
\[\left(\frac{d!}{2^d}\cdot n^d\cdot\ln\left(n^d\right)\right)^{-1/(d+1)}\cdot\tau_n^\mathsf{cc}\longrightarrow1\quad\text{in probability.}\]
\end{thm}

\begin{proof}[Proof of Theorem \ref{thm:couponcollector}]
The proof is based on a first and second moment argument.

\medskip

\emph{Upper bound.}
Fix a parameter $c\in\mathbb{R}_+^*$, and set (in hindsight)
\[k_n(c)=\left\lfloor\left(c\cdot n^d\cdot\ln\left(n^d\right)\right)^{1/(d+1)}\right\rfloor\quad\text{for all $n\in\mathbb{N}^*$.}\]
We want to show that for each $c>d!\cdot2^{-d}$, we have $\tau_n^\mathsf{cc}\leq k_n (c)$ with high probability as $n\to\infty$.
To this end, let us calculate the first moment of the random variables $\left.\mathbb{T}_n^d\middle\backslash\mathcal{B}_n^\mathsf{cc}(k)\right.$.
We have
\[\mathbb{E}\left[\#\left(\mathbb{T}_n^d\middle\backslash\mathcal{B}_n^\mathsf{cc}(k)\right)\right]=\sum_{x\in\mathbb{T}_n^d}\mathbb{P}\left(x\notin\mathcal{B}_n^\mathsf{cc}(k)\right),\]
where for every $x\in\mathbb{T}_n^d$,
\[\begin{split}
\mathbb{P}\left(x\notin\mathcal{B}_n^\mathsf{cc}(k)\right)&=\mathbb{P}\left(\bigcap_{i=1}^k\left(X_i\notin\overline{B}(x,k-i)\right)\right)\\
&=\prod_{i=1}^k\left(1-\frac{\#\overline{B}(x,k-i)}{n^d}\right)=\prod_{j=0}^{k-1}\left(1-\frac{\#\overline{B}(0,j)}{n^d}\right).
\end{split}\]
Now, we claim that as $n\to\infty$, we have
\begin{eqnarray}
\prod_{j=0}^{k_n(c)-1}\left(1-\frac{\#\overline{B}(0,j)}{n^d}\right)&=&\exp\left(-\sum_{j=0}^{k_n(c)-1}\frac{\#\overline{B}(0,j)}{n^d}+o(1)\right)\label{eq:couponcollector1stmoment1}\\
&=&\exp\left(-\frac{2^d}{(d+1)!}\cdot c\cdot\ln\left(n^d\right)\cdot(1+o(1))\right)\label{eq:couponcollector1stmoment2}
\end{eqnarray}
Taking this for granted, let us conclude the proof of the upper bound.
A naive first moment argument would consist in upper bounding the probability that $\mathbb{T}_n^d$ is not burned by time $k_n(c)$ as follows:
\[\begin{split}
\mathbb{P}\left(\tau_n^\mathsf{cc}>k_n(c)\right)&=\mathbb{P}\left(\left.\mathbb{T}_n^d\middle\backslash\mathcal{B}_n^\mathsf{cc}(k_n(c))\right.\neq\emptyset\right)\\
&\leq\mathbb{E}\left[\#\left(\mathbb{T}_n^d\middle\backslash\mathcal{B}_n^\mathsf{cc}(k_n(c))\right)\right]=n^d\cdot\prod_{j=0}^{k_n(c)-1}\left(1-\frac{\#\overline{B}(0,j)}{n^d}\right).
\end{split}\]
However, by \eqref{eq:couponcollector1stmoment2}, it would require $c>(d+1)!\cdot2^{-d}$ for the right-hand side to go to zero as $n\to\infty$, and this threshold is actually off by a factor of $(d+1)$.
To improve on this first moment argument, let $\left(\overline{B}(x,r_n)\,;\,x\in A_n\right)$ be a covering of $\mathbb{T}_n^d$ by balls of radius $r_n:=\left\lfloor n^{d/(d+1)}\right\rfloor$, with centres $x\in\mathbb{T}_n^d$ more than $r_n$ apart from each other.
Note that such a covering can be constructed iteratively as in the proof of Proposition \ref{prop:detbounds}.
Next, since the balls $\left(\overline{B}(x,\lfloor r_n/2\rfloor)\,;\,x\in A_n\right)$ are disjoint, we have ${\#A_n\leq n^d\cdot\left(\#\overline{B}(0,\lfloor r_n/2\rfloor)\right)^{-1}}$.
In particular, we get that $\#A_n\leq C\cdot n^{d/(d+1)}$ for some constant $C=C(d)\in(0,\infty)$ that does not depend on $n$.
Now, notice that we have the inclusion of events
\[\left(\tau_n^\mathsf{cc}>k_n(c)+r_n\right)=\left(\left.\mathbb{T}_n^d\middle\backslash\mathcal{B}_n^\mathsf{cc}(k_n(c)+r_n)\right.\neq\emptyset\right)\subset\bigcup_{x\in A_n}\left(x\notin\mathcal{B}_n^\mathsf{cc}(k_n(c))\right).\]
Thus, by the union bound, we obtain
\[\begin{split}
\mathbb{P}\left(\tau_n^\mathsf{cc}>k_n(c)+r_n\right)&\leq\sum_{x\in A_n}\mathbb{P}\left(x\notin\mathcal{B}_n^\mathsf{cc}(k_n(c))\right)\\
&=\#A_n\cdot\prod_{j=0}^{k_n(c)-1}\left(1-\frac{\#\overline{B}(0,j)}{n^d}\right)\leq C\cdot n^{d/(d+1)}\cdot\prod_{j=0}^{k_n(c)-1}\left(1-\frac{\#\overline{B}(0,j)}{n^d}\right).
\end{split}\]
Now, by \eqref{eq:couponcollector1stmoment2}, it only takes $c>d!\cdot2^{-d}$ to make the right-hand side go to zero as $n\to\infty$.
This yields the desired upper bound on $\tau_n^\mathsf{cc}$, since $k_n(c)+r_n\sim\left(c\cdot n^d\cdot\ln\left(n^d\right)\right)^{1/(d+1)}$ as $n\to\infty$.

Finally, it remains to prove \eqref{eq:couponcollector1stmoment1} and \eqref{eq:couponcollector1stmoment2}.
To this end, we write
\[\prod_{j=0}^{k_n(c)-1}\left(1-\frac{\#\overline{B}(0,j)}{n^d}\right)=\exp\left(-\sum_{j=0}^{k_n(c)-1}\frac{\#\overline{B}(0,j)}{n^d}+\Delta_n\right),\]
where
\[\Delta_n=\sum_{j=0}^{k_n(c)-1}\left(\frac{\#\overline{B}(0,j)}{n^d}+\ln\left(1-\frac{\#\overline{B}(0,j)}{n^d}\right)\right).\]
Then, since
\[\max_{j\in\llbracket0,k_n(c)-1\rrbracket}\frac{\#\overline{B}(0,j)}{n^d}\leq\frac{\#\overline{B}(0,k_n(c))}{n^d}\underset{n\to\infty}{=}n^{-d/(d+1)+o(1)}\longrightarrow0,\]
we may bound $\Delta_n$ by using the inequality $|u+\ln(1-u)|\leq u^2$ for all $u\in[0,1/2]$.
For all $n\in\mathbb{N}^*$ sufficiently large so that $\#\overline{B}(0,k_n(c))\cdot n^{-d}\leq1/2$, we get
\[|\Delta_n|\leq\sum_{j=0}^{k_n(c)-1}\left(\frac{\#\overline{B}(0,j)}{n^d}\right)^2\leq k_n(c)\cdot\left(\frac{\#\overline{B}(0,k_n(c))}{n^d}\right)^2\underset{n\to\infty}{=}n^{-d/(d+1)+o(1)}\longrightarrow0,\]
which proves \eqref{eq:couponcollector1stmoment1}.
Equation \eqref{eq:couponcollector1stmoment2} now readily follows: we have
\[\sum_{j=0}^{k_n(c)-1}\frac{\#\overline{B}(0,j)}{n^d}\sim\frac{1}{n^d}\cdot\frac{2^d}{d!}\cdot\frac{k_n(c)^{d+1}}{d+1}\sim\frac{2^d}{(d+1)!}\cdot c\cdot\ln\left(n^d\right)\quad\text{as $n\to\infty$.}\]

\medskip

\emph{Lower bound.}
Let us prove that for each $c<d!\cdot2^{-d}$, we have $\tau_n^\mathsf{cc}>k_n(c)$ with high probability as $n\to\infty$.
To this end, let us calculate the second moment of the random variable $\#\left(\mathbb{T}_n^d\middle\backslash\mathcal{B}_n^\mathsf{cc}(k_n(c))\right)$.
Indeed, by the Cauchy--Schwarz inequality, we have
\[\mathbb{P}\left(\tau_n^\mathsf{cc}>k_n(c)\right)=\mathbb{P}\left(\left.\mathbb{T}_n^d\middle\backslash\mathcal{B}_n^\mathsf{cc}(k_n(c))\right.\neq\emptyset\right)\geq\frac{\mathbb{E}\left[\#\left(\mathbb{T}_n^d\middle\backslash\mathcal{B}_n^\mathsf{cc}(k_n(c))\right)\right]^2}{\mathbb{E}\left[\#\left(\mathbb{T}_n^d\middle\backslash\mathcal{B}_n^\mathsf{cc}(k_n(c))\right)^2\right]}.\]
Thus, it suffices to show that for each $c<d!\cdot2^{-d}$, we have
\begin{equation}\label{eq:couponcollector2ndmomentgoal}
\frac{\mathbb{E}\left[\#\left(\mathbb{T}_n^d\middle\backslash\mathcal{B}_n^\mathsf{cc}(k_n(c))\right)^2\right]}{\mathbb{E}\left[\#\left(\mathbb{T}_n^d\middle\backslash\mathcal{B}_n^\mathsf{cc}(k_n(c))\right)\right]^2}\longrightarrow1\quad\text{as $n\to\infty$.}
\end{equation}
The second moment of the random variable $\#\left(\mathbb{T}_n^d\middle\backslash\mathcal{B}_n^\mathsf{cc}(k_n(c))\right)$ can be calculated as follows: we have
\[\mathbb{E}\left[\#\left(\mathbb{T}_n^d\middle\backslash\mathcal{B}_n^\mathsf{cc}(k_n(c))\right)^2\right]=\sum_{x_1,x_2\in\mathbb{T}_n^d}\mathbb{P}\left(x_1,x_2\notin\mathcal{B}_n^\mathsf{cc}(k_n(c))\right),\]
where for every $x_1,x_2\in\mathbb{T}_n^d$,
\[\begin{split}
\mathbb{P}\left(x_1,x_2\notin\mathcal{B}_n^\mathsf{cc}(k_n(c))\right)&=\mathbb{P}\left(\bigcap_{i=1}^{k_n(c)}\left(X_i\notin\left(\overline{B}(x_1,k_n(c)-i)\cup\overline{B}(x_2,k_n(c)-i)\right)\right)\right)\\
&=\prod_{i=1}^{k_n(c)}\left(1-\frac{\#\left(\overline{B}(x_1,k_n(c)-i)\cup\overline{B}(x_2,k_n(c)-i)\right)}{n^d}\right)\\
&=\prod_{j=0}^{k_n(c)-1}\left(1-\frac{\#\left(\overline{B}(0,j)\cup\overline{B}(x_2-x_1,j)\right)}{n^d}\right).
\end{split}\]
Thus, we get
\[\mathbb{E}\left[\#\left(\mathbb{T}_n^d\middle\backslash\mathcal{B}_n^\mathsf{cc}(k_n(c))\right)^2\right]=n^d\cdot\sum_{x\in\mathbb{T}_n^d}\prod_{j=0}^{k_n(c)-1}\left(1-\frac{\#\left(\overline{B}(0,j)\cup\overline{B}(x,j)\right)}{n^d}\right),\]
and it follows that
\begin{eqnarray*}
\lefteqn{\frac{\mathbb{E}\left[\#\left(\mathbb{T}_n^d\middle\backslash\mathcal{B}_n^\mathsf{cc}(k_n(c))\right)^2\right]}{\mathbb{E}\left[\#\left(\mathbb{T}_n^d\middle\backslash\mathcal{B}_n^\mathsf{cc}(k_n(c))\right)\right]^2}}\\
&=&\frac{1}{n^d}\cdot\sum_{x\in\mathbb{T}_n^d}\prod_{j=0}^{k_n(c)-1}\left(\left(1-\frac{\#\left(\overline{B}(0,j)\cup\overline{B}(x,j)\right)}{n^d}\right)\cdot\left(1-\frac{\#\overline{B}(0,j)}{n^d}\right)^{-2}\right).
\end{eqnarray*}
Now, we claim that as $n\to\infty$, we have
\begin{equation}\label{eq:couponcollector2ndmoment}
\prod_{j=0}^{k_n(c)-1}\left(1-\frac{\#\left(\overline{B}(0,j)\cup\overline{B}(x,j)\right)}{n^d}\right)=\exp\left(-\sum_{j=0}^{k_n(c)-1}\frac{\#\left(\overline{B}(0,j)\cup\overline{B}(x,j)\right)}{n^d}+o(1)\right)
\end{equation}
uniformly in $x\in\mathbb{T}_n^d$.
Indeed, we can write
\[\prod_{j=0}^{k_n(c)-1}\left(1-\frac{\#\left(\overline{B}(0,j)\cup\overline{B}(x,j)\right)}{n^d}\right)=\exp\left(-\sum_{j=0}^{k_n(c)-1}\frac{\#\left(\overline{B}(0,j)\cup\overline{B}(x,j)\right)}{n^d}+\Delta_n(x)\right),\]
where
\[\Delta_n(x)=\sum_{j=0}^{k_n(c)-1}\left(\frac{\#\left(\overline{B}(0,j)\cup\overline{B}(x,j)\right)}{n^d}+\ln\left(1-\frac{\#\left(\overline{B}(0,j)\cup\overline{B}(x,j)\right)}{n^d}\right)\right).\]
Then, since
\[\max_{j\in\llbracket0,k_n(c)-1\rrbracket}\frac{\#\left(\overline{B}(0,j)\cup\overline{B}(x,j)\right)}{n^d}\leq\frac{2\cdot\#\overline{B}(0,k_n(c))}{n^d}\underset{n\to\infty}{=}n^{-d/(d+1)+o(1)}\longrightarrow0,\]
we may bound $\Delta_n(x)$ by using the inequality $|u+\ln(1-u)|\leq u^2$ for all $u\in[0,1/2]$.
For all $n\in\mathbb{N}^*$ sufficiently large so that $\left(2\cdot\#\overline{B}(0,k_n(c))\right)\cdot n^{-d}\leq1/2$, we get
\[|\Delta_n(x)|\leq\sum_{j=0}^{k_n(c)-1}\left(\frac{\#\left(\overline{B}(0,j)\cup\overline{B}(x,j)\right)}{n^d}\right)^2\leq k_n(c)\cdot\left(\frac{2\cdot\#\overline{B}(0,k_n(c))}{n^d}\right)^2.\]
This proves \eqref{eq:couponcollector2ndmoment}, since the right-hand side does not depend on $x$ and goes to zero as $n\to\infty$.

Now, combining \eqref{eq:couponcollector1stmoment1} and \eqref{eq:couponcollector2ndmoment}, we get that as $n\to\infty$, we have
\begin{eqnarray*}
\lefteqn{\prod_{j=0}^{k_n(c)-1}\left(\left(1-\frac{\#\left(\overline{B}(0,j)\cup\overline{B}(x,j)\right)}{n^d}\right)\cdot\left(1-\frac{\#\overline{B}(0,j)}{n^d}\right)^{-2}\right)}\\
&=&\exp\left(\sum_{j=0}^{k_n(c)-1}\frac{\#\left(\overline{B}(0,j)\cap\overline{B}(x,j)\right)}{n^d}+o(1)\right)\quad\text{uniformly in $x\in\mathbb{T}_n^d$.}
\end{eqnarray*}
We deduce that
\begin{equation}\label{eq:couponcollector2ndmomenteq}
\frac{\mathbb{E}\left[\#\left(\mathbb{T}_n^d\middle\backslash\mathcal{B}_n^\mathsf{cc}(k_n(c))\right)^2\right]}{\mathbb{E}\left[\#\left(\mathbb{T}_n^d\middle\backslash\mathcal{B}_n^\mathsf{cc}(k_n(c))\right)\right]^2}\underset{n\to\infty}{\sim}\frac{1}{n^d}\cdot\sum_{x\in\mathbb{T}_n^d}\exp\left(\sum_{j=0}^{k_n(c)-1}\frac{\#\left(\overline{B}(0,j)\cap\overline{B}(x,j)\right)}{n^d}\right).
\end{equation}
Finally, we have everything at hand to prove that for each $c<d!\cdot2^{-d}$, the convergence \eqref{eq:couponcollector2ndmomentgoal} holds.
We upper bound the right-hand side of \eqref{eq:couponcollector2ndmomenteq} as follows: since $\overline{B}(0,j)\cap\overline{B}(x,j)=\emptyset$ as soon as $d(0,x)>2k_n(c)$, we have
\[\frac{1}{n^d}\cdot\sum_{x\in\mathbb{T}_n^d}\exp\left(\sum_{j=0}^{k_n(c)-1}\frac{\#\left(\overline{B}(0,j)\cap\overline{B}(x,j)\right)}{n^d}\right)\leq1+\frac{1}{n^d}\cdot\#\overline{B}(0,2k_n(c))\cdot\exp\left(\sum_{j=0}^{k_n(c)-1}\frac{\#\overline{B}(0,j)}{n^d}\right).\]
Now, observe that on the one hand, we have $\#\overline{B}(0,2k_n(c))\cdot n^{-d}=n^{-d/(d+1)+o(1)}$ as $n\to\infty$, and on the other hand, we have
\[\sum_{j=0}^{k_n(c)-1}\frac{\#\overline{B}(0,j)}{n^d}\sim\frac{1}{n^d}\cdot\frac{2^d}{d!}\cdot\frac{k_n(c)^{d+1}}{d+1}\sim\frac{2^d}{d!}\cdot c\cdot\ln\left(n^{d/(d+1)}\right)\quad\text{as $n\to\infty$,}\]
with $\left.2^d\middle/d!\right.\cdot c<1$ by assumption.
Therefore, we obtain that
\[\frac{1}{n^d}\cdot\#\overline{B}(0,2k_n(c))\cdot\exp\left(\sum_{j=0}^{k_n(c)-1}\frac{\#\overline{B}(0,j)}{n^d}\right)\longrightarrow0\quad\text{as $n\to\infty$,}\]
which yields \eqref{eq:couponcollector2ndmomentgoal}.
This completes the proof of the lower bound, and of the theorem.
\end{proof}

\end{document}
